\documentclass[12pt]{article}
\usepackage[utf8]{inputenc}
\usepackage[T1]{fontenc}
\usepackage{hyperref}
\title{Twana Language Resources for K-12 Curriculum Development}
\author{Compiled by Claude Sonnet 5 (claude-sonnet-5) for Andr\'as Kornai}
\date{}
\begin{document}
\maketitle

\section{Introduction}
Twana (Tuwaduq, ISO 639-3 twa) is a Central Coast Salish language of the Hood Canal area of Washington State. It is closely related to Lushootseed. The last fluent speaker reportedly died in 1980. The Skokomish Indian Tribe is now working on revival.

\section{Primary Educational Resources}
\begin{itemize}
\item The Twana Language Project, led by Christie Chambless, has reportedly taught vocabulary to about 200 middle-school students (as reported by ICT News). No published curriculum was found. \url{https://ictnews.org/?p=38979}
\end{itemize}

\section{Linguistic and Reference Materials}
\begin{itemize}
\item tuwaduq: The Twana Language E-Dictionary Project (UBC RelLex record). The Skokomish Indian Tribe is the responsible party, and the dictionary draws on earlier work by Drachman and Elmendorf and on recordings of former speakers. \url{https://wiki.ubc.ca/Documentation:RelLex/tuwaduq:_The_Twana_Language_E-Dictionary_Project}
\item Wikipedia, Twana language. \url{https://en.wikipedia.org/wiki/Twana_language}
\end{itemize}

\section{Religious and Community Resources}
\begin{itemize}
\item Wikipedia, Twana people. \url{https://en.wikipedia.org/wiki/Twana}
\end{itemize}
The Skokomish Indian Tribe is the community body responsible for the project. I did not see an official tribal URL in the search results.

\section{Digital and Audio Resources}
The online Twana dictionary (released 2020 according to Wikipedia) is described under Reference Materials. I did not see its direct URL in the results.

\section{Curriculum Development Context}
The language has no living first-language speakers. Revival rests on the tribal dictionary and project, and school use so far appears to be vocabulary teaching in middle school. Any curriculum should be developed with the Skokomish Indian Tribe.
\end{document}